\clearpage

%\section {Letter}

	\begin{center}
	
	%{\huge \textbf{Proposal Letter to  Local Emergency Planning Committee }}

	{\textbf{Proposal Letter to  Local Emergency Planning Committee }}

    
	\end{center}


%\noindent\rule{\textwidth}{1pt}
%\vspace{-0.2cm}
%\noindent \ To: 
%\vspace{-0.2cm}
%\noindent \ From: Team \#2515426 of 2025 MCM
%\vspace{-0.2cm}
%\noindent \ Date: February 31, 2025
%\vspace{-0.2cm}
%\noindent \ Subject:
%\vspace{-0.4cm}
%\noindent\rule{\textwidth}{1pt}


% 在需要背景图的页面开头添加
\backgroundsetup{
  scale=1,              % 图片缩放比例（1为原尺寸）
  angle=0,              % 图片旋转角度（0为不旋转）
  opacity=0.1,          % 透明度（0-1，0为完全透明）
  contents={\includegraphics[width=\paperwidth,height=\paperheight]{background.png}}  % 图片路径及尺寸
}
\BgThispage  % 仅当前页应用此背景


\leftline{Dear Chairman and Members of the Committee:}
%开头顶格左对齐
First, we would like to express our sincere admiration for the contributions of the Committee to the protection of public safety in the region. We are particularly grateful for your forward-looking work in emergency planning, which provids a valuable foundation for our research.

Our team has developed an innovative \textbf{"Dynamic Zonal Collaborative Search" }strategy based on mathematical modeling and simulation analysis. By abstracting building spaces into network structures, this method has not only demonstrated exceptional performance in simulating typical single-story office buildings but also verified its universality in complex environments such as multi-story buildings and non-standard layouts. This breakthrough is attributed to three core proprietary technologies:

Firstly, our "Non-Return Intelligent Path Planning System" revolutionizes the inefficient back-and-forth movement of responders in traditional searches through advanced ant colony optimization algorithms. It acts like a precise "building GPS" for each responder, automatically generating optimal continuous paths. For example, in multi-story buildings, the algorithm prioritizes deploying responders near stairwells to facilitate vertical coordination, ensuring a seamless journey from start to finish and guiding responders to the exit closest to the task completion point, thus completely eliminating time wasted on redundant movement.

Secondly, our proprietary "Risk-Aware Dynamic Prioritization System" enables optimal allocation of rescue resources. It calculates real-time risk coefficients for each area and prioritizes high-risk locations such as kindergartens. In multi-story buildings, it focuses on floors with dense population. Key results indicate that compared with traditional experience-dependent models, strategically deploying responders and prioritizing high-risk areas can reduce total search time by over 30\%, achieving a perfect balance between safety and speed within the precious rescue window.

Most importantly, our system boasts adaptability to complex disasters. The model can accurately predict the spread of thick smoke or toxic gas leaks in fires. The system can optimize paths and adjust team sizes through testing, ensuring that response strategies evolve intelligently with disaster development.

Based on the above technologies, we propose some suggestions to further enhance emergency management effectiveness:


1. \textbf{Integrate technologies to support dynamic decision-making}: Deploy IoT sensors to monitor smoke, temperature, and responder positions in real-time. Feed this data into the adaptive model to provide responders with immediate guidance when conditions change. Simultaneously link with building automation systems to prioritize alarm handling and optimize communication links, ensuring unobstructed information flow.

2. \textbf{Strengthen practical drills and training}: Conduct regular drills in simulated smoke environments to familiarize teams with model-optimized routes, focusing on high-risk areas and incorporating extreme scenarios such as communication failures and insufficient responders. This helps cultivate muscle memory for orderly evacuation, ensuring that algorithmic strategies can be accurately implemented in actual emergencies.

By combining scientific path planning, intelligent sensing technologies, and solid drills, emergency evacuation searches can be transformed from chaotic escapes into orderly transfers. We earnestly request the Committee to adopt these strategies to maximize the protection of lives. Our model and complete report are available for further discussion, and we look forward to working with you to elevate the region’s emergency management capabilities to new heights, building a stronger technological defense for public safety.

\rightline{Sincerely,}%注意这里要空一整行（敲两个回车），这样才会分段，一个回车和另个回车一样都默认是连在一起的

\rightline{Team 16793}
\rightline{Nov. 17,2025}


